\chapter*{Abstract}\label{abstract}
The rapid adoption of generative Artificial Intelligence (AI) is reshaping introductory programming education, providing new opportunities to scaffold learning while simultaneously creating a crisis for academic integrity. Because AI tools can now generate functional code instantly from natural language prompts, they have made traditional assessments highly vulnerable to cheating. Furthermore, unguided use of these tools without pedagogical guardrails or structured instruction leads to a fundamental over-reliance, where students may outsource the problem-solving process entirely and bypass the critical cognitive struggles necessary for conceptual mastery.

This thesis focuses on building and testing two specific AI-driven interventions designed to strategically integrate these technologies into the curriculum. First, a course-grounded AI tutor was developed using Retrieval-Augmented Generation (RAG) to provide students with pedagogical guidance strictly aligned with official course materials, ensuring that the AI acts as a tutor rather than a solution-generator. Second, a scalable AI-driven oral assessment system was implemented to conduct interactive interviews, probing a student’s conceptual understanding of their submitted work to ensure it reflects their own reasoning. These interventions were evaluated within the context of a pilot course at UNSW (COMP9021) using a mixed-methods approach. This evaluation combines quantitative interaction metrics, such as system usage patterns and performance data, with qualitative student surveys to assess the impact on learning outcomes, perceived fairness, and the effectiveness of the integration. By shifting the focus from the final code output to the authentic problem-solving process, this study seeks to enable the benefits of AI to scale while preserving the fundamental principles of computational thinking.

\noindent\textbf{Keywords:} Artificial Intelligence, Programming Education, Academic Integrity, Retrieval-Augmented Generation (RAG), Authentic Assessment, Mixed-Methods Evaluation